\chapter{Discussion}
\label{chap:discussion}

% Sources: docs/Design Principles.md, docs/Optimality Scope.md,
%          docs/Justification of Implementation.md §4

Chapter~\ref{chap:experiments} reported the results experiment by experiment. This chapter synthesizes them into five broader findings. Section~\ref{sec:disc:synthesis} discusses when and why scheduling improves performance, Section~\ref{sec:disc:scope} defines the scope of the resulting optimality claims, Section~\ref{sec:disc:optimistic} discusses the heralded-versus-optimistic timing trade-off exploited by the searched schedules, and Section~\ref{sec:disc:limitations} summarizes the limitations of the model and search procedure.

\section{When Scheduling Matters}
\label{sec:disc:synthesis}

The five findings below concern the \emph{performance and feasibility advantages} obtained by searching over purification schedules. They generalize beyond the reference configuration $\mathcal{N}_\mathrm{ref}$, subject to the search scope of Section~\ref{sec:disc:scope}.

\paragraph{Feasibility and rate advantages}
Under the idealized noise model tested by the source papers, the paper schedule and the searched schedules are all feasible, and the gap between them is a modest rate improvement ($2.5\%$ at matched cost, Section~\ref{sec:exp:headline}). Introducing inner-qubit error, which is not modelled in the source papers, can make the paper schedule fall below the fidelity floor. At the same generation cost, the optimizer's search can instead select a feasible schedule (Section~\ref{sec:exp:netgen}). In both cases, the mechanism is the same: the searched schedule \emph{concentrates purification where noise has accumulated most}, which a uniform hand-designed schedule cannot do by construction.

\paragraph{The scheduling advantage generalizes beyond hand-picked configurations}
A sample of 100 random per-hop heterogeneous networks confirms the same advantage beyond a single hand-chosen configuration: the searched schedule beats a uniform convention in 80 of 100 cases, by a mean of $20.37\%$ in rate, and never does worse (Section~\ref{sec:exp:heterogeneous}). One deliberately constructed network, with a single hop assigned a $12\times$ higher inner-qubit error rate than its neighbours, similarly causes the paper schedule to become infeasible while a feasible schedule can still be found by search, as detailed in Table~\ref{tab:weaklink} (Appendix~\ref{app:weaklink}).

\paragraph{The advantage concentrates in an intermediate regime}
\label{sec:disc:nearoptimal}
The noise sweep (Section~\ref{sec:exp:edsweep}) and the Pareto frontier (Section~\ref{sec:exp:pareto}) both show that scheduling matters most when the fidelity constraint becomes difficult to attain. On the other hand, scheduling barely matters when the fidelity floor is easy to satisfy, whether because noise is low or the allowed cost is generous. In that regime, most schedules are feasible. Their rates thus cluster together, and the searched frontier itself goes flat past a certain cost (already $C=60$ at the reference noise level). The rate improvement over the uniform baseline increases from near zero at $e_d=0$ to $18.2\%$ at $e_d=0.01$, with a clear difference emerging above about $e_d=0.005$. However, this advantage does not persist at higher noise levels. With $e_d=0.02$, both the uniform baseline and the budget-relaxed schedule fall below $F_{\min}=0.9$ once $e_d$ reaches approximately $0.014$--$0.016$. Beyond this point, neither schedule is feasible, so comparing their rates is no longer meaningful. Scheduling therefore matters most at intermediate noise levels, where the fidelity constraint is restrictive enough to distinguish between schedules but still allows feasible solutions.

In the noise sweep's fidelity panel (Figure~\ref{fig:edsweep:fidelity}), the paper baseline sits above both optimizer variants for nearly the whole range. This does \textbf{not} mean that it performs better than the optimizer: once $F_{\min}=0.9$ is reached, additional fidelity provides \textbf{no benefit}. The paper schedule always spends the full $C=100$ budget on a fixed circuit, while the budget-relaxed optimizer uses \emph{less}, starting from $C=20$ near $e_d=0$ and increasing towards $C=100$ as noise grows. It therefore accepts fidelity closer to the floor in exchange for higher success probability and rate, as shown in Figure~\ref{fig:edsweep:rate}. At higher noise, this trade-off also extends the feasible range: the fixed paper schedule falls below $F_{\min}=0.9$ by $e_d=0.014$, while both optimizer variants remain feasible by re-optimizing the schedule as $e_d$ increases.

\paragraph{The same constraint becomes harder to sustain with chain length}
The minimum generation cost required to achieve $F\ge0.9$ grows faster than the linear $10N$ budget proposed in the source paper over the tested range, with a descriptive fit of $E_{\max}^{\min}\approx1.19\,N^{1.67}$ (Section~\ref{sec:exp:budget}). The fitted curve first exceeds the $10N$ budget at $N=28$. This fit is descriptive rather than asymptotic, and the crossover is sensitive to the beam width, as discussed in Section~\ref{sec:alg:beam}. Even at $N=10$, the paper's proposed budget already exceeds the minimum required cost by a factor of $2.0$--$4.8$ across the tested noise range. These observations reflect the same underlying effect seen in the noise sweep: \emph{a fixed fidelity floor becomes harder to sustain as either the chain length or the noise increases}. Section~\ref{sec:disc:scope} revisits the budget-scaling result below, where limitations of the searched families become important when interpreting this scaling.

\paragraph{The search accepts other objectives without requiring new algorithms}
None of the results above needed a different search procedure. Replacing the objective, rate, fidelity, or cost changes only which schedule is returned, not how the search finds it (Section~\ref{sec:exp:objectives}), because candidate construction and evaluation are kept separate throughout Chapter~\ref{chap:algorithms}.

\section{Optimality Scope}
\label{sec:disc:scope}

% Source: docs/Optimality Scope.md, Optimizer Status.md §Known Limits

The beam-search tier (Section~\ref{sec:alg:beam}) includes one same-span pumping operation, with the candidate frontier bounded by the beam width $w_{\mathrm{beam}}$. All returned schedules are therefore valid constructions within the searched families, not proofs of global optimality. Finding the same best candidate when the search is run with two different beam widths provides evidence that the result is stable with respect to this parameter, but does not prove that no better candidate was pruned at either width. An uncapped beam width pumping mode (Section~\ref{sec:alg:pumping}) is available but becomes impractical on larger instances, so it is used only for validation on smaller networks rather than as the default search method.

The \emph{limited expressiveness} of the search space is demonstrated by the following example. The optimizer can only construct schedules from the primitive blocks and schedule families introduced in Chapter~\ref{chap:algorithms}, and therefore cannot represent every structurally valid schedule. For example, it cannot construct a schedule by combining two independently constructed full-span candidates. At $N=18$, such a manually constructed schedule fits within the Integrating paper's budget of $E_{\max}=10N=180$, while the optimizer finds a schedule requiring an even smaller minimum budget (Section~\ref{sec:exp:budget}). The two results are therefore not in conflict: the latter concerns only schedules within the optimizer's search space. More generally, the budget-scaling result gives the minimum budget among the searched families, not a lower bound on the budget required by all legal schedules.

The results of Chapter~\ref{chap:experiments} should therefore distinguish between existence and non-existence claims. An existence claim, that a feasible schedule was found, holds for the searched model because the schedule is \emph{explicitly constructed}, validated, and evaluated by the shared physical model. A non-existence claim applies only to the families and search limits actually explored and does not prove that no legal schedule exists outside that space.

\section{Heralded vs.\ Optimistic Execution}
\label{sec:disc:optimistic}

% Source: Validated Formal Model Def.md §2.6, sweep_gamma_and_tau_emit/README.md

The timing results in Section~\ref{sec:exp:timing} show that the relative benefit of optimistic and heralded execution depends on both \emph{memory dephasing} and \emph{generation latency}. The description of herald placement in Section~\ref{sec:model:dag} captures this timing trade-off. An optimistic sequence defers heralding until later in the sequence, eliminating intermediate classical round-trips but allowing subsequent operations to proceed before the outcome of an earlier purification is known. At the reference configuration, with no inner-qubit error and $\gamma=0$, the optimizer \emph{consistently favours optimistic candidates} because the reduced communication delay \textbf{outweighs} the loss in success probability. When $\gamma>0$, heralded pumping is further penalized because its sacrificial copies remain in memory during each round-trip delay, increasing their exposure to decoherence.

Under the tested conditions, these effects converge in the same direction: reducing communication latency also reduces the time spent waiting in memory and therefore limits additional decoherence. The trade-off depends on the communication delay, memory coherence, and failure probabilities. With different physical parameters, intermediate heralding could instead be preferable. The observed preference should therefore be understood as a result of the configurations tested here, rather than as a general architectural rule.

\section{Limitations}
\label{sec:disc:limitations}

% Source: Optimizer Status.md §Known Limits

The results are subject to limitations in the schedule space, physical model, and bounded search.

\begin{enumerate}
    \item \textit{Non-adaptive schedules.} Only non-adaptive schedules are evaluated: the schedule is fixed before execution and cannot condition a later purification step on an intermediate outcome (Section~\ref{sec:model:scope}). A fully adaptive strategy would require a different formulation, so the results bound only the non-adaptive schedule space.

    \item \textit{Memory footprint.} The search does not enforce limits on memory footprint. $M_{\max}$, the number of branches held in memory at once, is computed exactly for every candidate (Section~\ref{sec:model:scope}), but only the generation-cost budget $E_{\max}$ constrains the search. A schedule feasible under $E_{\max}$ may therefore still require more simultaneously held branches than a deployment with a strict memory limit could support.

    \item \textit{Noise model.} The fidelity model assumes Pauli-channel noise, consistent with the source papers. Schedule comparisons are therefore internally consistent, but their relative performance could change under hardware dominated by different error processes. In particular, coupled loss, gate-error, and memory effects are not represented.

    \item \textit{Fixed beam width.} The search is sensitive to which construction types compete for retained slots. Enabling same-span pumping, for example, changes which candidate is returned at the reference configuration, so adding a construction can change previously reported results even at fixed beam width. Extensions should therefore include explicit before-and-after comparisons rather than assume that a new construction can only help.
\end{enumerate}